\section{Static Operating Payoffs and Optimal Controls}

This phase derives the benchmark static operating blocks exactly as they are written in the benchmark manuscript. The goal is to obtain the optimized flow payoffs that will later enter the Bellman equations. We do not redesign the model, we do not expand argument lists, and we do not insert post-optimization objects into primitive objectives.

\subsection{Benchmark Fidelity Check}

We first restate the exact benchmark objects that this phase must derive.

\paragraph{Type-$A$ benchmark object.}
The benchmark manuscript writes
\[
\pi_A(z;w,M)
=
\pi_A^G(z;w)+\pi_A^{O,\mathrm{inc}}(z;w,M)-f_A
+x_A\Omega_A(z)-\frac{\chi_A}{\eta}x_A^\eta,
\qquad \eta>1,
\]
with
\[
\Omega_A(z)\equiv \lambda_A(z)R_A(z).
\]

\paragraph{Type-$B$ benchmark object.}
The benchmark manuscript writes
\[
H_B(z;p_m,M)=\lambda_B(z,p_m,M)\bigl(R_B(z)-p_m\bigr)-\tau_M(M),
\]
and then
\[
\pi_B(z;p_m,M)
=
-f_B+x_BH_B(z;p_m,M)-\frac{\chi_B}{\eta}x_B^\eta.
\]
So the primitive type-$B$ objective uses $x_BH_B$, not $x_B[H_B]_+$. The positive-part operator must be derived from the optimization problem.

\paragraph{Type-$C$ benchmark object.}
The benchmark manuscript writes
\[
\pi_C(z;p_m,w,M)
=
\pi_C^G(z;w)+\pi_C^{O,\mathrm{inc}}(z;w,M)-f_C
+\max_{m\ge 0}\{p_m m-c(m,z;w)\}.
\]

\paragraph{State space and compact parameter domains used in this phase.}
From Phase 1, the capability state lives in a compact metric space $Z\subset \mathbb R_+$. We now restate that compactness because we will use it directly in continuity and boundedness claims. Whenever we discuss continuity or boundedness with respect to prices, we fix compact price domains
\[
\mathcal P_m\subset \mathbb R_+,
\qquad
\mathcal W\subset \mathbb R_+.
\]
These compact sets are not new model objects. They are local parameter domains used only to state regularity results carefully. We will use the shorthand
\[
K_A:=Z,\qquad K_B:=Z\times \mathcal P_m,\qquad K_C:=Z\times \mathcal P_m\times \mathcal W.
\]

\paragraph{Scope of this phase.}
This phase derives
\[
x_A^*(z),\quad \Pi_A(z;w,M),\quad
x_B^*(z;p_m,M),\quad \Pi_B(z;p_m,M),\quad
m^*(z;p_m,w),\quad \Pi_C(z;p_m,w,M).
\]

\subsection{Type-$A$ firms}

\paragraph{Economic role.}
Type $A$ firms are integrated firms. In the benchmark they already produce medicines and also choose original-drug idea-generation intensity. The control variable is $x_A$. The object $\Pi_A(z;w,M)$ is the optimized static payoff that will later appear inside the Bellman equations.

\paragraph{Step 1: write the exact static problem.}
Fix $(z,w,M)$. The benchmark static problem is
\begin{equation}
\max_{x_A\ge 0}
\left\{
\pi_A^G(z;w)+\pi_A^{O,\mathrm{inc}}(z;w,M)-f_A
+x_A\Omega_A(z)-\frac{\chi_A}{\eta}x_A^\eta
\right\}.
\label{eq:p2A_problem_rev}
\end{equation}
The first three terms do not depend on $x_A$. Therefore, if we define the reduced objective
\[
g_A(x_A;z):=x_A\Omega_A(z)-\frac{\chi_A}{\eta}x_A^\eta,
\]
then maximizing the benchmark payoff in \eqref{eq:p2A_problem_rev} is equivalent to maximizing $g_A(x_A;z)$ over $x_A\ge 0$.

\paragraph{Step 2: state the choice variable, state, and given objects.}
The choice variable is $x_A\in[0,\infty)$. The state variable is $z\in Z$. The objects taken as given are $w$, $M$, $\Omega_A(z)$, $\chi_A$, and $\eta$.

\paragraph{Step 3: assumptions used.}
\textbf{Benchmark assumptions:} $\eta>1$, $\chi_A>0$, and $\Omega_A(z)=\lambda_A(z)R_A(z)$.  
\textbf{Supplemental regularity assumptions:}
\begin{enumerate}
    \item[(A1)] For each $z\in Z$, the number $\Omega_A(z)$ is finite.
    \item[(A2)] For the benchmark closed form without a positive-part operator, $\Omega_A(z)\ge 0$ for all $z\in Z$.
    \item[(A3)] The functions $z\mapsto \pi_A^G(z;w)$, $z\mapsto \pi_A^{O,\mathrm{inc}}(z;w,M)$, and $z\mapsto \Omega_A(z)$ are continuous on the compact set $Z$.
\end{enumerate}

\paragraph{Step 4: prove existence of an optimum.}
We want to apply the Weierstrass Extreme Value Theorem. That theorem requires a continuous objective on a nonempty compact feasible set. Our original feasible set $[0,\infty)$ is not compact, so we first show that the optimizer must lie in a bounded interval.

First, by A1 the map $x_A\mapsto g_A(x_A;z)$ is continuous on $[0,\infty)$ for every fixed $z$. Second, because $\eta>1$,
\[
g_A(x_A;z)
=
x_A\Omega_A(z)-\frac{\chi_A}{\eta}x_A^\eta
=
x_A\left[\Omega_A(z)-\frac{\chi_A}{\eta}x_A^{\eta-1}\right].
\]
As $x_A\to\infty$, the term $x_A^{\eta-1}$ diverges to $+\infty$, so the bracketed term tends to $-\infty$. Multiplying by the positive term $x_A$ gives
\[
g_A(x_A;z)\to -\infty
\qquad\text{as }x_A\to\infty.
\]
Therefore there exists a number $R_A(z)>0$ such that
\[
g_A(x_A;z)<g_A(0;z)=0
\qquad\text{for all }x_A\ge R_A(z).
\]
So no maximizer can lie to the right of $R_A(z)$. Hence any maximizer must belong to the interval $[0,R_A(z)]$.

Now the Weierstrass theorem applies because:
\begin{enumerate}
    \item the feasible set $[0,R_A(z)]$ is nonempty;
    \item the feasible set $[0,R_A(z)]$ is compact in $\mathbb R$;
    \item the function $g_A(\cdot;z)$ is continuous on that compact set.
\end{enumerate}
So $g_A(\cdot;z)$ attains a maximum on $[0,R_A(z)]$. Therefore the benchmark type-$A$ problem has at least one solution.

\paragraph{Step 5: compute the first derivative.}
For $x_A>0$,
\[
\frac{\partial}{\partial x_A}\left(x_A\Omega_A(z)\right)=\Omega_A(z),
\qquad
\frac{\partial}{\partial x_A}\left(\frac{\chi_A}{\eta}x_A^\eta\right)=\chi_A x_A^{\eta-1}.
\]
Therefore
\begin{equation}
g_A'(x_A;z)=\Omega_A(z)-\chi_A x_A^{\eta-1}.
\label{eq:p2A_derivative_rev}
\end{equation}
Because $\eta>1$, the right-hand side extends continuously to $x_A=0$ by setting $0^{\eta-1}=0$, so
\[
g_A'(0;z)=\Omega_A(z).
\]

\paragraph{Step 6: compute the second derivative and explain sufficiency.}
For $x_A>0$,
\begin{equation}
g_A''(x_A;z)=-\chi_A(\eta-1)x_A^{\eta-2}.
\label{eq:p2A_second_derivative_rev}
\end{equation}
Since $\chi_A>0$, $\eta-1>0$, and $x_A^{\eta-2}>0$ for $x_A>0$, we have
\[
g_A''(x_A;z)<0
\qquad\text{for all }x_A>0.
\]
So $g_A$ is strictly concave on $(0,\infty)$. Equivalently, because $x_A\mapsto x_A^\eta$ is strictly convex when $\eta>1$, the function $x_A\mapsto -(\chi_A/\eta)x_A^\eta$ is strictly concave, and adding the linear term $x_A\Omega_A(z)$ preserves strict concavity. Therefore any feasible point satisfying the Kuhn--Tucker conditions is the unique global maximizer.

\paragraph{Step 7: derive the policy rule carefully.}
We now solve the problem by cases.

\textit{Case 1: $\Omega_A(z)=0$.}
In this case,
\[
g_A'(x_A;z)=0-\chi_A x_A^{\eta-1}=-\chi_A x_A^{\eta-1}<0
\qquad\text{for every }x_A>0.
\]
So the reduced objective is strictly decreasing on $(0,\infty)$. Therefore the largest feasible value is attained at the boundary:
\[
x_A^*(z)=0.
\]

\textit{Case 2: $\Omega_A(z)>0$.}
In this case,
\[
g_A'(0;z)=\Omega_A(z)>0.
\]
From \eqref{eq:p2A_second_derivative_rev}, the derivative $g_A'(\cdot;z)$ is strictly decreasing on $(0,\infty)$. Also,
\[
g_A'(x_A;z)=\Omega_A(z)-\chi_A x_A^{\eta-1}\to -\infty
\qquad\text{as }x_A\to\infty.
\]
So $g_A'$ starts positive at $0$, then decreases continuously, and eventually becomes negative. By the Intermediate Value Theorem, there exists a unique $x_A>0$ such that
\[
g_A'(x_A;z)=0.
\]
That unique root solves
\[
\Omega_A(z)-\chi_A \bigl(x_A^*(z)\bigr)^{\eta-1}=0.
\]
Rearranging gives
\[
\chi_A \bigl(x_A^*(z)\bigr)^{\eta-1}=\Omega_A(z),
\]
and dividing by $\chi_A$ gives
\[
\bigl(x_A^*(z)\bigr)^{\eta-1}=\frac{\Omega_A(z)}{\chi_A}.
\]
Taking both sides to the power $1/(\eta-1)$ yields
\[
x_A^*(z)=\left(\frac{\Omega_A(z)}{\chi_A}\right)^{\frac{1}{\eta-1}}.
\]

\textit{Combine the two cases.}
Because A2 implies $\Omega_A(z)\ge 0$, the two cases may be combined as
\begin{equation}
x_A^*(z)=\left(\frac{\Omega_A(z)}{\chi_A}\right)^{\frac{1}{\eta-1}}.
\label{eq:p2A_policy_rev}
\end{equation}

\paragraph{Step 8: substitute the policy back into the payoff step by step.}
The optimized benchmark payoff is
\[
\Pi_A(z;w,M)
=
\pi_A^G(z;w)+\pi_A^{O,\mathrm{inc}}(z;w,M)-f_A
+x_A^*(z)\Omega_A(z)-\frac{\chi_A}{\eta}\bigl(x_A^*(z)\bigr)^\eta.
\]
From the first-order condition in the interior case,
\[
\chi_A \bigl(x_A^*(z)\bigr)^{\eta-1}=\Omega_A(z).
\]
Multiply both sides by $x_A^*(z)$:
\[
\chi_A \bigl(x_A^*(z)\bigr)^\eta=x_A^*(z)\Omega_A(z).
\]
Now divide both sides by $\eta$:
\[
\frac{\chi_A}{\eta}\bigl(x_A^*(z)\bigr)^\eta=\frac{1}{\eta}x_A^*(z)\Omega_A(z).
\]
Substitute this into the payoff:
\begin{align*}
\Pi_A(z;w,M)
&=
\pi_A^G(z;w)+\pi_A^{O,\mathrm{inc}}(z;w,M)-f_A
+x_A^*(z)\Omega_A(z)-\frac{1}{\eta}x_A^*(z)\Omega_A(z)\\
&=
\pi_A^G(z;w)+\pi_A^{O,\mathrm{inc}}(z;w,M)-f_A
+\left(1-\frac{1}{\eta}\right)x_A^*(z)\Omega_A(z)\\
&=
\pi_A^G(z;w)+\pi_A^{O,\mathrm{inc}}(z;w,M)-f_A
+\frac{\eta-1}{\eta}x_A^*(z)\Omega_A(z).
\end{align*}
Next substitute the closed-form expression for $x_A^*(z)$:
\[
x_A^*(z)\Omega_A(z)
=
\left(\frac{\Omega_A(z)}{\chi_A}\right)^{\frac{1}{\eta-1}}\Omega_A(z).
\]
Now separate the numerator and denominator powers:
\[
\left(\frac{\Omega_A(z)}{\chi_A}\right)^{\frac{1}{\eta-1}}
=
\Omega_A(z)^{\frac{1}{\eta-1}}\chi_A^{-\frac{1}{\eta-1}}.
\]
Therefore
\begin{align*}
x_A^*(z)\Omega_A(z)
&=
\Omega_A(z)^{\frac{1}{\eta-1}}\chi_A^{-\frac{1}{\eta-1}}\Omega_A(z)^1\\
&=
\chi_A^{-\frac{1}{\eta-1}}\Omega_A(z)^{\frac{1}{\eta-1}+1}.
\end{align*}
Since
\[
\frac{1}{\eta-1}+1=\frac{1}{\eta-1}+\frac{\eta-1}{\eta-1}=\frac{\eta}{\eta-1},
\]
we obtain
\[
x_A^*(z)\Omega_A(z)
=
\chi_A^{-\frac{1}{\eta-1}}\Omega_A(z)^{\frac{\eta}{\eta-1}}.
\]
Hence
\begin{equation}
\Pi_A(z;w,M)
=
\pi_A^G(z;w)+\pi_A^{O,\mathrm{inc}}(z;w,M)-f_A
+\frac{\eta-1}{\eta}\chi_A^{-\frac{1}{\eta-1}}\Omega_A(z)^{\frac{\eta}{\eta-1}}.
\label{eq:p2A_payoff_rev}
\end{equation}

\paragraph{Step 9: continuity and boundedness of the derived objects.}
Because $Z$ is compact and $\Omega_A$ is continuous on $Z$ by A3, the Extreme Value Theorem implies that $\Omega_A$ is bounded on $Z$. Let
\[
\bar\Omega_A:=\sup_{z\in Z}\Omega_A(z)<\infty.
\]
Define the function
\[
h_A(u):=\left(\frac{u}{\chi_A}\right)^{\frac{1}{\eta-1}}
\qquad\text{for }u\in[0,\bar\Omega_A].
\]
The map $u\mapsto u/\chi_A$ is continuous, and the map $v\mapsto v^{1/(\eta-1)}$ is continuous on $[0,\infty)$. Therefore $h_A$ is continuous on $[0,\bar\Omega_A]$. Since $x_A^*(z)=h_A(\Omega_A(z))$, the policy function $z\mapsto x_A^*(z)$ is continuous on $Z$ as a composition of continuous functions.

Next define
\[
q_A(u):=\frac{\eta-1}{\eta}\chi_A^{-\frac{1}{\eta-1}}u^{\frac{\eta}{\eta-1}}
\qquad\text{for }u\in[0,\bar\Omega_A].
\]
By the same continuity argument, $q_A$ is continuous on $[0,\bar\Omega_A]$. Equation \eqref{eq:p2A_payoff_rev} may be written as
\[
\Pi_A(z;w,M)=\pi_A^G(z;w)+\pi_A^{O,\mathrm{inc}}(z;w,M)-f_A+q_A(\Omega_A(z)).
\]
Every term on the right-hand side is continuous in $z$, so $\Pi_A(\cdot;w,M)$ is continuous on $Z$. Since $Z$ is compact, continuity implies boundedness. Thus both $x_A^*$ and $\Pi_A$ are continuous and bounded on $K_A=Z$.

\subsubsection*{What did we just do, and why does it matter?}
We solved the integrated firm's one-period research problem exactly as it appears in the benchmark manuscript. The economic logic is that type $A$ raises research intensity up to the point where the marginal value of idea generation equals the marginal convex cost. Because the objective is strictly concave, this balancing condition determines a unique optimum. The resulting object $\Pi_A(z;w,M)$ is the flow payoff that later lets the dynamic model compare operating as an integrated firm with the other organizational forms.

\paragraph{What has been established mathematically.}
We proved existence and uniqueness for the benchmark type-$A$ static problem, derived the policy rule $x_A^*(z)$, and obtained the optimized payoff $\Pi_A(z;w,M)$.

\paragraph{What assumptions were used.}
We used $\eta>1$, $\chi_A>0$, the benchmark definition $\Omega_A(z)=\lambda_A(z)R_A(z)$, the nonnegativity condition $\Omega_A(z)\ge 0$, and continuity of the primitive objects on compact $Z$.

\paragraph{What remains to be derived next.}
We next derive the benchmark type-$B$ static problem, paying particular attention to the fact that $[H_B]_+$ is a derived object rather than a primitive object.

\subsection{Type-$B$ firms}

\paragraph{Economic role.}
Type $B$ firms are research-specialized firms. In the benchmark they do not produce internally. Their control variable is the original-drug idea-generation intensity $x_B$. The object $H_B(z;p_m,M)$ is the one-idea net commercialization value. The optimized payoff $\Pi_B(z;p_m,M)$ is the static object that later enters the Bellman equations for type $B$.

\paragraph{Step 1: write the exact static problem.}
Fix $(z,p_m,M)$. The benchmark first defines
\begin{equation}
H_B(z;p_m,M)=\lambda_B(z,p_m,M)\bigl(R_B(z)-p_m\bigr)-\tau_M(M).
\label{eq:p2B_HB_rev}
\end{equation}
The exact current payoff is
\begin{equation}
\pi_B(z;p_m,M)
=
-f_B+x_BH_B(z;p_m,M)-\frac{\chi_B}{\eta}x_B^\eta.
\label{eq:p2B_payoff_primitive_rev}
\end{equation}
So the benchmark static problem is
\begin{equation}
\max_{x_B\ge 0}
\left\{
-f_B+x_BH_B(z;p_m,M)-\frac{\chi_B}{\eta}x_B^\eta
\right\}.
\label{eq:p2B_problem_rev}
\end{equation}
Since $-f_B$ does not depend on $x_B$, we may equivalently maximize
\[
g_B(x_B;z,p_m,M):=x_BH_B(z;p_m,M)-\frac{\chi_B}{\eta}x_B^\eta
\]
over $x_B\ge 0$.

\paragraph{Step 2: state the choice variable, state, and given objects.}
The choice variable is $x_B\in[0,\infty)$. The state is $z\in Z$. The objects taken as given are $p_m\in\mathcal P_m$, $M$, $H_B(z;p_m,M)$, $\chi_B$, and $\eta$.

\paragraph{Step 3: assumptions used.}
\textbf{Benchmark assumptions:} $\eta>1$, $\chi_B>0$, and the benchmark definition \eqref{eq:p2B_HB_rev}.  
\textbf{Supplemental regularity assumptions:}
\begin{enumerate}
    \item[(B1)] For each $(z,p_m)\in K_B$, the number $H_B(z;p_m,M)$ is finite.
    \item[(B2)] The map $(z,p_m)\mapsto H_B(z;p_m,M)$ is continuous on the compact set $K_B=Z\times\mathcal P_m$.
\end{enumerate}

\paragraph{Step 4: prove existence of an optimum.}
As in the type-$A$ problem, we want to use the Weierstrass theorem, so we first reduce the noncompact feasible set to a compact interval.

For fixed $(z,p_m,M)$, the map $x_B\mapsto g_B(x_B;z,p_m,M)$ is continuous on $[0,\infty)$ because $H_B(z;p_m,M)$ is finite by B1. Also,
\[
g_B(x_B;z,p_m,M)
=
x_B\left[H_B(z;p_m,M)-\frac{\chi_B}{\eta}x_B^{\eta-1}\right].
\]
Since $\eta>1$, the term $x_B^{\eta-1}$ diverges to $+\infty$ as $x_B\to\infty$, so the bracketed term tends to $-\infty$. Therefore
\[
g_B(x_B;z,p_m,M)\to -\infty
\qquad\text{as }x_B\to\infty.
\]
Hence there exists a number $R_B(z,p_m,M)>0$ such that
\[
g_B(x_B;z,p_m,M)<g_B(0;z,p_m,M)=0
\qquad\text{for all }x_B\ge R_B(z,p_m,M).
\]
So any maximizer must lie in the compact interval $[0,R_B(z,p_m,M)]$.

Now the Weierstrass theorem applies because:
\begin{enumerate}
    \item the feasible set $[0,R_B(z,p_m,M)]$ is nonempty and compact;
    \item the function $g_B(\cdot;z,p_m,M)$ is continuous on that set.
\end{enumerate}
Therefore the benchmark type-$B$ problem has at least one solution.

\paragraph{Step 5: compute the first derivative.}
For $x_B>0$,
\[
\frac{\partial}{\partial x_B}\left(x_BH_B(z;p_m,M)\right)=H_B(z;p_m,M),
\qquad
\frac{\partial}{\partial x_B}\left(\frac{\chi_B}{\eta}x_B^\eta\right)=\chi_B x_B^{\eta-1}.
\]
Hence
\begin{equation}
g_B'(x_B;z,p_m,M)=H_B(z;p_m,M)-\chi_B x_B^{\eta-1}.
\label{eq:p2B_derivative_rev}
\end{equation}
As before, because $\eta>1$, this derivative extends continuously to $x_B=0$ with value
\[
g_B'(0;z,p_m,M)=H_B(z;p_m,M).
\]

\paragraph{Step 6: compute the second derivative and explain sufficiency.}
For $x_B>0$,
\begin{equation}
g_B''(x_B;z,p_m,M)=-\chi_B(\eta-1)x_B^{\eta-2}.
\label{eq:p2B_second_derivative_rev}
\end{equation}
Since $\chi_B>0$, $\eta-1>0$, and $x_B^{\eta-2}>0$ for $x_B>0$, we have
\[
g_B''(x_B;z,p_m,M)<0
\qquad\text{for all }x_B>0.
\]
Thus $g_B$ is strictly concave on $(0,\infty)$, and indeed on all of $[0,\infty)$ because the nonlinear cost term is strictly convex. Therefore any feasible point satisfying the Kuhn--Tucker conditions is the unique global maximizer.

\paragraph{Step 7: derive the policy rule and show where $[H_B]_+$ comes from.}
This is the benchmark-critical step. We start from the primitive objective containing $H_B$, not $[H_B]_+$.

\textit{Case 1: $H_B(z;p_m,M)\le 0$.}
Take any $x_B>0$. Then
\[
g_B'(x_B;z,p_m,M)=H_B(z;p_m,M)-\chi_B x_B^{\eta-1}\le 0-\chi_B x_B^{\eta-1}<0.
\]
So for every positive $x_B$, the derivative is strictly negative. That means the objective is strictly decreasing on $(0,\infty)$. Therefore the largest feasible value is attained at the boundary:
\[
x_B^*(z;p_m,M)=0.
\]

\textit{Case 2: $H_B(z;p_m,M)>0$.}
In this case,
\[
g_B'(0;z,p_m,M)=H_B(z;p_m,M)>0.
\]
From \eqref{eq:p2B_second_derivative_rev}, the derivative $g_B'(\cdot;z,p_m,M)$ is strictly decreasing on $(0,\infty)$. Also,
\[
g_B'(x_B;z,p_m,M)=H_B(z;p_m,M)-\chi_B x_B^{\eta-1}\to -\infty
\qquad\text{as }x_B\to\infty.
\]
So $g_B'$ starts positive at zero, decreases continuously, and eventually becomes negative. By the Intermediate Value Theorem, there exists a unique $x_B>0$ such that
\[
g_B'(x_B;z,p_m,M)=0.
\]
That unique root solves
\[
H_B(z;p_m,M)-\chi_B \bigl(x_B^*(z;p_m,M)\bigr)^{\eta-1}=0.
\]
Rearrange:
\[
\chi_B \bigl(x_B^*(z;p_m,M)\bigr)^{\eta-1}=H_B(z;p_m,M),
\]
divide by $\chi_B$:
\[
\bigl(x_B^*(z;p_m,M)\bigr)^{\eta-1}=\frac{H_B(z;p_m,M)}{\chi_B},
\]
and raise both sides to the power $1/(\eta-1)$:
\[
x_B^*(z;p_m,M)=\left(\frac{H_B(z;p_m,M)}{\chi_B}\right)^{\frac{1}{\eta-1}}.
\]

\textit{Compact representation.}
Define the positive-part operator by
\[
[u]_+:=\max\{u,0\}.
\]
If $H_B(z;p_m,M)\le 0$, then $[H_B(z;p_m,M)]_+=0$ and the policy is $x_B^*=0$. If $H_B(z;p_m,M)>0$, then $[H_B(z;p_m,M)]_+=H_B(z;p_m,M)$ and the policy is the interior formula above. Therefore the two cases can be summarized exactly as
\begin{equation}
x_B^*(z;p_m,M)=\left(\frac{[H_B(z;p_m,M)]_+}{\chi_B}\right)^{\frac{1}{\eta-1}}.
\label{eq:p2B_policy_rev}
\end{equation}
This is the correct place for $[H_B]_+$ to appear. It is not primitive; it is a compact summary of the solved optimization problem.

\paragraph{Step 8: substitute the policy back into the payoff step by step.}
We again proceed by cases.

\textit{Case 1: $H_B(z;p_m,M)\le 0$.}
Then $x_B^*(z;p_m,M)=0$, so
\[
\Pi_B(z;p_m,M)=-f_B.
\]

\textit{Case 2: $H_B(z;p_m,M)>0$.}
Then
\[
\Pi_B(z;p_m,M)
=
-f_B+x_B^*H_B(z;p_m,M)-\frac{\chi_B}{\eta}(x_B^*)^\eta.
\]
From the first-order condition,
\[
\chi_B(x_B^*)^{\eta-1}=H_B.
\]
Multiply both sides by $x_B^*$:
\[
\chi_B(x_B^*)^\eta=x_B^*H_B.
\]
Divide both sides by $\eta$:
\[
\frac{\chi_B}{\eta}(x_B^*)^\eta=\frac{1}{\eta}x_B^*H_B.
\]
Substitute this identity into the payoff:
\begin{align*}
\Pi_B(z;p_m,M)
&=
-f_B+x_B^*H_B-\frac{1}{\eta}x_B^*H_B\\
&=
-f_B+\left(1-\frac{1}{\eta}\right)x_B^*H_B\\
&=
-f_B+\frac{\eta-1}{\eta}x_B^*H_B.
\end{align*}
Now substitute the interior closed-form policy:
\[
x_B^*H_B
=
\left(\frac{H_B}{\chi_B}\right)^{\frac{1}{\eta-1}}H_B.
\]
Separate the powers:
\[
\left(\frac{H_B}{\chi_B}\right)^{\frac{1}{\eta-1}}
=
H_B^{\frac{1}{\eta-1}}\chi_B^{-\frac{1}{\eta-1}}.
\]
Therefore
\begin{align*}
x_B^*H_B
&=
H_B^{\frac{1}{\eta-1}}\chi_B^{-\frac{1}{\eta-1}}H_B^1\\
&=
\chi_B^{-\frac{1}{\eta-1}}H_B^{\frac{1}{\eta-1}+1}.
\end{align*}
Since
\[
\frac{1}{\eta-1}+1=\frac{\eta}{\eta-1},
\]
we obtain
\[
x_B^*H_B=\chi_B^{-\frac{1}{\eta-1}}H_B^{\frac{\eta}{\eta-1}}.
\]
So in the case $H_B>0$,
\[
\Pi_B(z;p_m,M)
=
-f_B+\frac{\eta-1}{\eta}\chi_B^{-\frac{1}{\eta-1}}H_B(z;p_m,M)^{\frac{\eta}{\eta-1}}.
\]
Combining the two cases with the positive-part operator gives
\begin{equation}
\Pi_B(z;p_m,M)
=
-f_B+\frac{\eta-1}{\eta}\chi_B^{-\frac{1}{\eta-1}}[H_B(z;p_m,M)]_+^{\frac{\eta}{\eta-1}}.
\label{eq:p2B_optimized_rev}
\end{equation}

\paragraph{Step 9: continuity and boundedness of the derived objects.}
Because $K_B=Z\times\mathcal P_m$ is compact and $H_B$ is continuous on $K_B$ by B2, the Extreme Value Theorem implies that $H_B$ is bounded on $K_B$. Let
\[
\bar H_B:=\sup_{(z,p_m)\in K_B}|H_B(z;p_m,M)|<\infty.
\]
Define
\[
h_B(u):=\left(\frac{[u]_+}{\chi_B}\right)^{\frac{1}{\eta-1}}
\qquad\text{for }u\in[-\bar H_B,\bar H_B].
\]
The positive-part map $u\mapsto [u]_+$ is continuous on $\mathbb R$. The map $v\mapsto v/\chi_B$ is continuous, and the map $r\mapsto r^{1/(\eta-1)}$ is continuous on $[0,\infty)$. Therefore $h_B$ is continuous on $[-\bar H_B,\bar H_B]$. Since
\[
x_B^*(z;p_m,M)=h_B(H_B(z;p_m,M)),
\]
the policy function $(z,p_m)\mapsto x_B^*(z;p_m,M)$ is continuous on $K_B$.

Next define
\[
q_B(u):=\frac{\eta-1}{\eta}\chi_B^{-\frac{1}{\eta-1}}[u]_+^{\frac{\eta}{\eta-1}}
\qquad\text{for }u\in[-\bar H_B,\bar H_B].
\]
Again, $q_B$ is continuous. Equation \eqref{eq:p2B_optimized_rev} may be written as
\[
\Pi_B(z;p_m,M)=-f_B+q_B(H_B(z;p_m,M)).
\]
So $\Pi_B$ is continuous on $K_B$ as a composition of continuous functions. Since $K_B$ is compact, continuity implies boundedness. Thus both $x_B^*$ and $\Pi_B$ are continuous and bounded on $K_B$.

\subsubsection*{What did we just do, and why does it matter?}
We solved the benchmark research-specialist problem without altering the primitive objective. This matters because the sign of $H_B(z;p_m,M)$ is the channel through which MAH changes the attractiveness of research specialization. If commercialization conditions are bad so that $H_B\le 0$, the firm shuts down idea generation and chooses zero effort. If commercialization conditions are good so that $H_B>0$, the firm chooses positive effort according to the same marginal-value-equals-marginal-cost logic as in type $A$. That is why $[H_B]_+$ belongs in the solved policy rule and optimized payoff, not in the primitive objective.

\paragraph{What has been established mathematically.}
We proved existence and uniqueness for the benchmark type-$B$ static problem, derived the policy rule, and obtained the optimized payoff $\Pi_B(z;p_m,M)$.

\paragraph{What assumptions were used.}
We used $\eta>1$, $\chi_B>0$, the benchmark definition of $H_B$, and the continuity and finiteness assumptions stated on the compact domain $K_B$.

\paragraph{What remains to be derived next.}
We next derive the benchmark manufacturing-specialist problem for type $C$.

\subsection{Type-$C$ firms}

\paragraph{Economic role.}
Type $C$ firms are manufacturing-specialized firms. In the benchmark they do not originate new original-drug ideas. Their control variable is qualified contract-manufacturing output $m$. The optimized payoff $\Pi_C(z;p_m,w,M)$ is the current flow payoff from operating as a manufacturing specialist.

\paragraph{Step 1: write the exact static problem.}
Fix $(z,p_m,w,M)$. The benchmark manuscript writes
\begin{equation}
\pi_C(z;p_m,w,M)
=
\pi_C^G(z;w)+\pi_C^{O,\mathrm{inc}}(z;w,M)-f_C
+\max_{m\ge 0}\left\{p_m m-c(m,z;w)\right\}.
\label{eq:p2C_benchmark_payoff_rev}
\end{equation}
So the static control problem is
\begin{equation}
\max_{m\ge 0}\left\{p_m m-c(m,z;w)\right\}.
\label{eq:p2C_problem_rev}
\end{equation}
For later use, define the reduced objective
\[
g_C(m;z,p_m,w):=p_m m-c(m,z;w).
\]
Once this problem is solved, the optimized benchmark payoff is
\begin{equation}
\Pi_C(z;p_m,w,M)
=
\pi_C^G(z;w)+\pi_C^{O,\mathrm{inc}}(z;w,M)-f_C
+p_m m^*(z;p_m,w)-c(m^*(z;p_m,w),z;w).
\label{eq:p2C_payoff_rev}
\end{equation}

\paragraph{Step 2: state the choice variable, state, and given objects.}
The choice variable is $m\in[0,\infty)$. The state is $z\in Z$. The objects taken as given are $p_m\in\mathcal P_m$, $w\in\mathcal W$, $M$, and the manufacturing cost function $c(m,z;w)$.

\paragraph{Step 3: assumptions used.}
\textbf{Benchmark assumption:} the payoff specification \eqref{eq:p2C_benchmark_payoff_rev}.  
\textbf{Supplemental regularity assumptions:}
\begin{enumerate}
    \item[(C1)] The function $(m,z,w)\mapsto c(m,z;w)$ is jointly continuous on $[0,\infty)\times Z\times \mathcal W$.
    \item[(C2)] For each fixed $(z,w)$, the map $m\mapsto c(m,z;w)$ is twice continuously differentiable and strictly convex on $[0,\infty)$.
    \item[(C3)] Uniform coercivity on $Z\times\mathcal W$:
    \[
    \lim_{m\to\infty}\ \inf_{(z,w)\in Z\times\mathcal W}\frac{c(m,z;w)}{m}=\infty.
    \]
    \item[(C4)] The partial derivative $c_m(m,z;w)$ is jointly continuous in $(m,z,w)$ on $[0,\infty)\times Z\times\mathcal W$.
\end{enumerate}

\paragraph{Step 4: prove existence of an optimum.}
We again want to use the Weierstrass theorem, so we first show that the relevant maximizers lie in a common compact interval.

Let
\[
\bar p:=\sup_{p_m\in\mathcal P_m} p_m.
\]
Because $\mathcal P_m$ is compact, $\bar p<\infty$. By C3, there exists a number $R_C>0$ such that for all $m\ge R_C$,
\[
\inf_{(z,w)\in Z\times\mathcal W}\frac{c(m,z;w)}{m}>\bar p+1.
\]
This implies that for every $(z,w)\in Z\times\mathcal W$ and every $m\ge R_C$,
\[
\frac{c(m,z;w)}{m}>\bar p+1.
\]
Multiplying by $m>0$ yields
\[
c(m,z;w)>(\bar p+1)m.
\]
Now fix any $(z,p_m,w)\in K_C$. Since $p_m\le \bar p$, we have for all $m\ge R_C$,
\begin{align*}
g_C(m;z,p_m,w)
&=
p_m m-c(m,z;w)\\
&<
p_m m-(\bar p+1)m\\
&\le
\bar p\,m-(\bar p+1)m\\
&=
-m.
\end{align*}
Therefore
\[
g_C(m;z,p_m,w)\to -\infty
\qquad\text{uniformly over }(z,p_m,w)\in K_C.
\]
In particular, for each fixed parameter triple, no maximizer can lie to the right of $R_C$. So every maximizer lies in the compact interval $[0,R_C]$.

Now the Weierstrass theorem applies to each fixed $(z,p_m,w)\in K_C$ because:
\begin{enumerate}
    \item the feasible set $[0,R_C]$ is nonempty and compact;
    \item by C1, the function $m\mapsto g_C(m;z,p_m,w)$ is continuous on $[0,R_C]$.
\end{enumerate}
Hence the type-$C$ problem has at least one solution for every parameter triple in $K_C$.

\paragraph{Step 5: compute the first derivative.}
We now define the notation
\[
c_m(m,z;w):=\frac{\partial c(m,z;w)}{\partial m}.
\]
For interior points $m>0$,
\[
\frac{\partial}{\partial m}\left(p_m m\right)=p_m,
\qquad
\frac{\partial}{\partial m}\left(c(m,z;w)\right)=c_m(m,z;w).
\]
Therefore
\begin{equation}
g_C'(m;z,p_m,w)=p_m-c_m(m,z;w).
\label{eq:p2C_derivative_rev}
\end{equation}

\paragraph{Step 6: compute the second derivative and explain sufficiency.}
Define
\[
c_{mm}(m,z;w):=\frac{\partial^2 c(m,z;w)}{\partial m^2}.
\]
Because C2 assumes twice continuous differentiability, the second derivative exists, and
\begin{equation}
g_C''(m;z,p_m,w)=-c_{mm}(m,z;w).
\label{eq:p2C_second_derivative_rev}
\end{equation}
Since $c(\cdot,z;w)$ is convex, we have $c_{mm}(m,z;w)\ge 0$ wherever the second derivative is evaluated, so $g_C''(m;z,p_m,w)\le 0$. This proves concavity.

For uniqueness, we need more than weak concavity: we use strict convexity of $c(\cdot,z;w)$. If $c(\cdot,z;w)$ is strictly convex, then $-c(\cdot,z;w)$ is strictly concave. Adding the linear term $m\mapsto p_m m$ preserves strict concavity. Therefore $g_C(\cdot;z,p_m,w)$ is strictly concave on $[0,\infty)$, and any feasible point satisfying the Kuhn--Tucker conditions is the unique global maximizer.

\paragraph{Step 7: derive the Kuhn--Tucker characterization carefully.}
Because the feasible set is $[0,\infty)$, the Kuhn--Tucker conditions are
\begin{equation}
p_m-c_m(m^*,z;w)\le 0,\qquad
m^*(z;p_m,w)\ge 0,\qquad
m^*(z;p_m,w)\bigl[p_m-c_m(m^*,z;w)\bigr]=0.
\label{eq:p2C_KKT_rev}
\end{equation}
We now analyze these conditions by cases.

\textit{Case 1: $p_m\le c_m(0,z;w)$.}
Take any $m>0$. Since $c(\cdot,z;w)$ is convex, its derivative $c_m(\cdot,z;w)$ is weakly increasing. Therefore
\[
c_m(m,z;w)\ge c_m(0,z;w)\ge p_m.
\]
Hence
\[
g_C'(m;z,p_m,w)=p_m-c_m(m,z;w)\le 0
\qquad\text{for every }m>0.
\]
So the objective is weakly decreasing on $(0,\infty)$, and the largest feasible value is attained at the boundary:
\begin{equation}
m^*(z;p_m,w)=0.
\label{eq:p2C_boundary_rev}
\end{equation}

\textit{Case 2: $p_m>c_m(0,z;w)$.}
Then
\[
g_C'(0;z,p_m,w)=p_m-c_m(0,z;w)>0.
\]
Because $g_C$ is differentiable at $0$, a positive derivative at zero means there exists $\delta>0$ such that
\[
g_C(\delta;z,p_m,w)>g_C(0;z,p_m,w).
\]
So the boundary point $m=0$ cannot be optimal. Since an optimum exists, the unique optimum must be an interior point. Any interior optimum must satisfy the first-order condition
\begin{equation}
c_m(m^*(z;p_m,w),z;w)=p_m.
\label{eq:p2C_interior_rev}
\end{equation}
Because $c(\cdot,z;w)$ is strictly convex, its derivative $c_m(\cdot,z;w)$ is strictly increasing. Therefore equation \eqref{eq:p2C_interior_rev} can have at most one solution. Since the optimizer exists and must be interior in this case, that solution is the unique optimizer.

\paragraph{Step 8: summarize the policy rule.}
Combining the two cases, the exact benchmark policy rule is
\begin{equation}
m^*(z;p_m,w)=
\begin{cases}
0, & \text{if } p_m\le c_m(0,z;w),\\[0.4em]
\text{the unique }m>0\text{ such that }c_m(m,z;w)=p_m, & \text{if } p_m>c_m(0,z;w).
\end{cases}
\label{eq:p2C_policy_rev}
\end{equation}
This is the correct benchmark characterization. A further closed form would require a special functional form for $c$, and the benchmark does not impose one here.

\paragraph{Step 9: state the optimized payoff carefully.}
Substituting the optimizer into the benchmark payoff gives
\begin{equation}
\Pi_C(z;p_m,w,M)
=
\pi_C^G(z;w)+\pi_C^{O,\mathrm{inc}}(z;w,M)-f_C
+p_m m^*(z;p_m,w)-c(m^*(z;p_m,w),z;w).
\label{eq:p2C_optimized_payoff_rev}
\end{equation}
There is no further benchmark-general algebraic simplification because no special functional form for $c$ has been imposed.

\paragraph{Step 10: prove continuity and boundedness carefully.}
We now invoke Berge's Maximum Theorem, so we state the assumptions we need explicitly.

Let the feasible correspondence be
\[
\Gamma(z,p_m,w):=[0,R_C]
\qquad\text{for every }(z,p_m,w)\in K_C.
\]
This correspondence is:
\begin{enumerate}
    \item nonempty, because $0\in[0,R_C]$;
    \item compact-valued, because $[0,R_C]$ is compact in $\mathbb R$;
    \item continuous as a correspondence, because it is constant.
\end{enumerate}
Next, by C1 the objective
\[
(m,z,p_m,w)\mapsto g_C(m;z,p_m,w)=p_m m-c(m,z;w)
\]
is jointly continuous on $[0,R_C]\times K_C$, which is compact. Therefore all assumptions of Berge's Maximum Theorem are satisfied: a jointly continuous objective and a nonempty compact-valued continuous feasible correspondence on a compact parameter domain.

Hence the value function
\[
V_C^{\mathrm{stat}}(z,p_m,w):=\max_{m\in[0,R_C]} g_C(m;z,p_m,w)
\]
is continuous on $K_C$. Because all maximizers of the original problem lie in $[0,R_C]$, this is the same object as
\[
V_C^{\mathrm{stat}}(z,p_m,w)=\max_{m\ge 0}\{p_m m-c(m,z;w)\}.
\]
Therefore the benchmark manufacturing block's optimized surplus term is continuous on $K_C$.

To obtain continuity of the policy rule, we use the standard unique-maximizer corollary to Berge's theorem: if the objective and feasible correspondence satisfy Berge's assumptions and the argmax is single-valued at every parameter point, then the argmax function is continuous. We have already established single-valuedness from strict concavity. Therefore
\[
(z,p_m,w)\mapsto m^*(z;p_m,w)
\]
is continuous on $K_C$.

Finally, because $K_C$ is compact and both $V_C^{\mathrm{stat}}$ and the primitive profit terms $\pi_C^G(z;w)$ and $\pi_C^{O,\mathrm{inc}}(z;w,M)$ are continuous on $K_C$, each of these objects is bounded on $K_C$ by the Extreme Value Theorem. Equation \eqref{eq:p2C_optimized_payoff_rev} then implies that $\Pi_C(z;p_m,w,M)$ is continuous and bounded on $K_C$.

\subsubsection*{What did we just do, and why does it matter?}
We solved the manufacturing specialist's one-period problem exactly as the benchmark writes it. Economically, type $C$ compares the contract-manufacturing price $p_m$ with marginal cost. If the price is weakly below marginal cost at zero output, it supplies nothing. If the price is above marginal cost at zero output, then the boundary cannot be optimal, and strict concavity implies a unique interior quantity where price equals marginal cost. The resulting object $\Pi_C(z;p_m,w,M)$ is the flow payoff that the dynamic model later uses to compare manufacturing specialization with the other organizational forms.

\paragraph{What has been established mathematically.}
We proved existence and uniqueness for the benchmark type-$C$ static problem, derived the Kuhn--Tucker characterization of $m^*(z;p_m,w)$, and obtained the optimized payoff $\Pi_C(z;p_m,w,M)$.

\paragraph{What assumptions were used.}
We used the benchmark type-$C$ payoff specification together with joint continuity, twice differentiability in $m$, strict convexity, and uniform coercivity on the compact domain $K_C$.

\paragraph{What remains to be derived next.}
All three benchmark static operating blocks are now derived. The next phase is to place $\Pi_A(z;w,M)$, $\Pi_B(z;p_m,M)$, and $\Pi_C(z;p_m,w,M)$ into the Bellman system and prove the relevant dynamic properties.

\subsection{Phase-2 output summary}

This phase has derived the exact benchmark optimized flow payoffs that enter the dynamic system:
\[
\Pi_A(z;w,M),\qquad \Pi_B(z;p_m,M),\qquad \Pi_C(z;p_m,w,M).
\]
For type $A$, we derived the unique research-intensity policy rule and its optimized payoff. For type $B$, we derived the policy rule from the primitive benchmark problem and showed why the positive-part operator $[H_B]_+$ is a derived compact notation rather than a primitive ingredient of the objective. For type $C$, we derived the unique contract-manufacturing supply rule through Kuhn--Tucker analysis and obtained the exact optimized payoff that appears in the benchmark.

These three objects are now ready to be inserted into the Bellman equations without any additional static rewriting.

\subsection{Assumptions used in this phase}

The derivations used the benchmark payoff specifications together with supplemental regularity conditions: continuity of primitive objects on compact domains, convex innovation costs for types $A$ and $B$, and a continuous, strictly convex, uniformly coercive manufacturing cost for type $C$.

\subsection{What remains for next phase}

Phase 3 will prove Bellman operator properties on $\mathcal C_b(Z)^3$:
\begin{enumerate}
    \item mapping from bounded-continuous triples into bounded-continuous triples,
    \item contraction under the sup norm,
    \item existence and uniqueness of $(V_A,V_B,V_C)$,
    \item existence of optimal organizational policy rules defined using the benchmark choice objects.
\end{enumerate}
